\begin{table}[H]
\centering
\caption{Data extraction form.}
\label{tab:DEF}
\resizebox{0.8\textwidth}{!}{
\begin{tabular}{|>{\centering\arraybackslash}p{0.1\textwidth}|>{\raggedright\arraybackslash}p{0.9\textwidth}|}
\hline
\rowcolor{gray!20}
 \raisebox{0.2\height}{\rule{0pt}{3ex}\textbf{RQs}} & \raisebox{0.2\height}{\rule{0pt}{3ex}\textbf{Evidence to be extracted}} \\
\hline
%\multicolumn{2}{|c|}{\textbf{Publication Space RQs}}\\  
%\hline

%PS-RQ1 & \begin{tabular}[c]{@{}l@{}}Publication year \\Publication type (article, conference paper, etc.)\end{tabular}\\
%\hline
%PS-RQ2 & \begin{tabular}[c]{@{}l@{}}Publication venue (conference name or journal name) \\Areas of interest of the venue (for journals we refer to the list of subject area \\ from Scimago; for conferences we refer to the source details from Scopus)\end{tabular}\\
%\hline
%PS-RQ3 & List of the authors’ affiliations \\
%\hline
%\multicolumn{2}{|c|}{\textbf{Research Space RQs}}\\  
%\hline
\RQlabel{RQ1} & List of sentences describing the adopted definitions of “Customer Churn” \\
\hline
\RQlabel{RQ2} & List of sentences identifying data/variables/features and the most effective features used for Churn Prediction \\
\hline
\RQlabel{RQ3} & List of sentences identifying data aggregation or correlation, or combination methods (e.g., based on time periods, order total value, products’ categories, seasonality, customer segmentation, etc.) that have been used before applying AI techniques on the resulting features \\
\hline
\RQlabel{RQ4} & List of sentences identifying the AI techniques used and found to be most effective for Churn Prediction in the Retail domain \\
\hline
\RQlabel{RQ5} & List of sentences identifying the metrics used to evaluate the performance of the adopted AI techniques. \\
\hline
\RQlabel{RQ6} & All the information related to the dataset used for Churn Prediction and publicly available (publication date, URL, \# of records, \# of customers, fields included, etc.) \\
\hline
\RQlabel{RQ7} & List of sentences on identified and discussed challenges and limitations \\
\hline
\RQlabel{RQ8} & List of sentences on emerging trends and future research directions \\
\cline{1-2}
\end{tabular}
}
\end{table}